\subsection{Deep Learning para dirección de arribo y beamforming}
\label{sec:dl_doa}

La estimación de dirección de arribo (DoA) es un problema central en radar, comunicaciones inalámbricas y procesamiento acústico.  
Los algoritmos clásicos como MUSIC destacan por su alta resolución angular, pero presentan limitaciones conocidas: baja robustez con pocas muestras (\emph{snapshots}), sensibilidad al ruido y dificultad para localizar fuentes coherentes.  
Estas restricciones reducen su eficacia en escenarios prácticos.

Un enfoque reciente, denominado \textbf{DLMUSIC}, combina técnicas de \emph{Deep Learning} con el marco clásico de MUSIC \cite{BaiChen2023_DLMUSIC}.  
La propuesta introduce dos innovaciones principales:
\begin{itemize}
	\item \textbf{Estimación de la covarianza:} en lugar de usar la matriz de covarianza empírica, se entrenan redes convolucionales (CNNs) para predecir una versión más estable y robusta, menos sensible al ruido y al número de snapshots.  
	\item \textbf{Espectro objetivo suavizado:} en vez de aproximar el espectro MUSIC real (que puede presentar picos distorsionados y lóbulos espurios en baja SNR), se utiliza un \emph{espectro idealizado} compuesto por funciones Gaussianas centradas en las direcciones verdaderas de llegada.  
	Esta representación suavizada actúa como objetivo de entrenamiento, guiando a la red hacia picos definidos y reduciendo falsos máximos.
\end{itemize}

Tras la estimación de la covarianza por las CNNs, el algoritmo procede con la descomposición en autovalores y la separación de subespacios propia de MUSIC, preservando así la interpretabilidad del método clásico.  
De esta manera, DLMUSIC no reemplaza la teoría de subespacios, sino que corrige las etapas más frágiles mediante aprendizaje profundo.

De acuerdo con los resultados reportados, este enfoque supera al MUSIC clásico en escenarios de \textbf{baja SNR}, con \textbf{pocos snapshots} y en presencia de \textbf{fuentes coherentes} \cite{BaiChen2023_DLMUSIC}.  
Ello demuestra cómo el \emph{Deep Learning} puede complementar a los algoritmos tradicionales, reforzando su desempeño en condiciones adversas y acercándolos a un uso más confiable en aplicaciones modernas.



\subsection{Relevancia y aplicaciones}

La importancia de enfoques como \textbf{DLMUSIC} radica en que resuelven varias de las limitaciones prácticas 
de MUSIC clásico.  
En escenarios reales, las señales rara vez cumplen las condiciones ideales: 
se dispone de pocas muestras, el ruido no es despreciable y las fuentes pueden estar correlacionadas.  
En esos casos, MUSIC pierde resolución o incluso falla en detectar las direcciones correctas.  

El aporte de DLMUSIC es reemplazar las etapas más frágiles —estimación de covarianza y espectro de referencia— 
por versiones aprendidas y suavizadas mediante \emph{deep learning}.  
Esto no elimina la interpretabilidad del método subespacial, pero lo hace más \textbf{robusto frente a ruido, 
	coherencia y baja cantidad de datos}.  
En la práctica, significa pasar de un algoritmo que funciona bien solo en condiciones de laboratorio 
a una herramienta con proyección en campo abierto.

Más allá de los mosquitos, esta combinación de aprendizaje profundo con teoría de subespacios 
puede aplicarse en dominios menos tradicionales, por ejemplo:
\begin{itemize}
	\item \textbf{Ubicación de fugas de gas} en plantas industriales usando arreglos de micrófonos para detectar ultrasonidos.  
	\item \textbf{Localización de jugadores} en estadios mediante arreglos acústicos para análisis deportivo en tiempo real.  
	\item \textbf{Identificación de puntos calientes en líneas de transmisión eléctrica} combinando antenas de RF y sensores térmicos.  
	\item \textbf{Seguimiento de fauna submarina} (peces o crustáceos) mediante redes de hidrófonos de bajo costo.  
\end{itemize}

Así, DLMUSIC no sólo refuerza la robustez de MUSIC en condiciones adversas, 
sino que abre la puerta a su uso en aplicaciones emergentes donde antes la estimación de dirección 
no era viable por falta de SNR o de datos.
